% Pastikan Anda menambahkan baris berikut di bagian preamble (sebelum \begin{document}):
% \usepackage{longtable}

% ========================
% DAFTAR SINGKATAN DAN ISTILAH
% ========================
\chapter*{DAFTAR SINGKATAN DAN ISTILAH}
\addcontentsline{toc}{chapter}{DAFTAR SINGKATAN DAN ISTILAH}

\vspace{0.5cm}

% Menggunakan 3 kolom: Istilah/Singkatan, Kepanjangan (Asing), dan Penjelasan
\setlength{\tabcolsep}{4pt}
\begin{longtable}{@{} L{2.4cm} L{3.0cm} L{7.4cm} @{}}
\textbf{Istilah / Singkatan} & \textbf{Kepanjangan / Asing} & \textbf{Penjelasan} \\[2ex]
\endhead % Memastikan header diulang ke halaman berikutnya

ABAC & Attribute-Based Access Control & Model kontrol akses yang memberi atau menolak izin berdasarkan evaluasi atribut subjek, objek, aksi, dan lingkungan. \\
\hline
ABAC Lab & --- & Platform dan repositori dataset ground-truth ABAC \parencite{bui_abac_2025} yang generator lognya diperluas pada penelitian ini. \\
\hline
ACL & Access Control List & Daftar kontrol akses; representasi izin per pasangan subjek-objek. \\
\hline
ACO & Ant Colony Optimization & Algoritma metaheuristik berbasis koloni semut; salah satu baseline berpopulasi \parencite{shang_abac_2024}. \\
\hline
Ancestral sampling & sampling ansestral & Pembangkitan solusi mengikuti urutan induk-anak pada struktur pohon dependensi. \\
\hline
Attribute-value pair & pasangan atribut-nilai & Unit penyusun aturan ABAC; banyaknya seluruh pasangan yang mungkin membentuk dimensi ruang solusi $D$. \\
\hline
Baseline & algoritma pembanding & Algoritma acuan yang dibandingkan terhadap algoritma usulan. \\
\hline
BIC & Bayesian Information Criterion & Kriteria pemilihan model yang dipakai BOA/iBOA saat mempelajari struktur dependensi. \\
\hline
BOA & Bayesian Optimization Algorithm & EDA multivariat yang membangun jaringan Bayes tanpa batasan orde dependensi. \\
\hline
CoDA-EDA & Compact Dependency-Aware Estimation of Distribution Algorithm & Algoritma yang dirancang dan diuji pada penelitian ini. \\
\hline
Chow-Liu tree & pohon Chow-Liu / pohon dependensi & Pohon rentang maksimum yang menghubungkan pasangan variabel ber-mutual information tertinggi (dibatasi orde satu). \\
\hline
compact-GA & compact Genetic Algorithm & Varian algoritma genetika hemat memori tanpa penyimpanan populasi eksplisit. \\
\hline
compact-PSO & compact Particle Swarm Optimization & Varian PSO hemat memori berbasis Probability Vector, diadaptasi ke ruang diskret. \\
\hline
Coverage & cakupan & Proporsi permintaan akses yang tercakup oleh sekurang-kurangnya satu aturan kebijakan. \\
\hline
$D$ & dimensi ruang solusi & Banyaknya pasangan atribut-nilai (beserta aksi) yang membentuk vektor keputusan biner. \\
\hline
DAC & Discretionary Access Control & Model kontrol akses diskresioner (dibahas pada tinjauan pustaka). \\
\hline
Drift & pergeseran distribusi data & Perubahan distribusi log akses antar-waktu yang memicu pembaruan struktur dependensi. \\
\hline
DSR & Design Science Research & Paradigma penelitian yang berorientasi pada perancangan dan evaluasi artefak. \\
\hline
DSRM & Design Science Research Methodology & Metodologi enam aktivitas Design Science Research \parencite{peffers_design_2007}. \\
\hline
Edge / edge computing & komputasi tepi & Pemrosesan data di dekat sumbernya, pada perangkat di tepi jaringan. \\
\hline
Edge-IoT & --- & Perangkat/gerbang tepi pada jaringan Internet of Things yang menjadi konteks target penelitian. \\
\hline
EDA & Estimation of Distribution Algorithm & Kelas algoritma evolusioner yang membangun dan men-sampling model probabilistik eksplisit dari solusi terbaik. \\
\hline
F1 / F-score & F-measure & Rata-rata harmonik presisi dan recall; metrik kualitas kebijakan utama. \\
\hline
Finite-sample bias & bias sampel-hingga & Bias estimator informasi (NMI/JSD) yang membesar seiring rasio jumlah sel kontingensi terhadap jumlah sampel. \\
\hline
FN & False Negative & Permintaan yang sebenarnya diizinkan namun diklasifikasikan sebagai tolak. \\
\hline
FP & False Positive & Permintaan yang sebenarnya ditolak namun diklasifikasikan sebagai izin (over-permission). \\
\hline
Friedman test & uji Friedman & Uji omnibus non-parametrik untuk mendeteksi perbedaan di antara banyak algoritma sekaligus. \\
\hline
GA & Genetic Algorithm & Algoritma genetika. \\
\hline
Gateway & gerbang & Perangkat penghubung di tepi jaringan IoT yang menjadi contoh lingkungan penerapan. \\
\hline
Ground-truth & kebijakan/keputusan acuan & Kebijakan atau label keputusan sebenarnya yang menjadi pembanding evaluasi. \\
\hline
H1--H4 & Hipotesis 1--4 & Empat hipotesis penelitian: kualitas kebijakan, jejak memori, efisiensi waktu pembaruan, dan warm-start. \\
\hline
Holm-Bonferroni & koreksi Holm-Bonferroni & Koreksi perbandingan berganda untuk mengendalikan inflasi galat Tipe I pada uji post-hoc. \\
\hline
iBOA & incremental Bayesian Optimization Algorithm & Varian BOA inkremental berbasis pohon dependensi (Chow-Liu). \\
\hline
IoT & Internet of Things & Jaringan perangkat fisik yang saling terhubung dan bertukar data. \\
\hline
JSD & Jensen-Shannon Divergence & Ukuran perbedaan dua distribusi probabilitas; dipakai untuk mengukur dan mendeteksi drift. \\
\hline
LLM & Large Language Model & Model bahasa besar; disebut sebagai arah riset di luar cakupan penelitian ini. \\
\hline
MAC & Mandatory Access Control & Model kontrol akses wajib (dibahas pada tinjauan pustaka). \\
\hline
MI & Mutual Information & Informasi bersama antar-variabel; dasar penyusunan pohon Chow-Liu. \\
\hline
Min-support & dukungan minimum & Ambang untuk menyatukan nilai atribut bersupport rendah (mis. pengenal) pada dataset besar. \\
\hline
NMI & Normalized Mutual Information & Mutual information ternormalisasi; ukuran korelasi atribut yang disuntikkan terkontrol ke dataset. \\
\hline
Noise floor & lantai derau & Nilai divergensi antar-jendela pada kondisi stasioner; menentukan kelayakan pemicu drift (H3). \\
\hline
Non-parametrik & non-parametric & Uji statistik yang tidak mengandaikan distribusi data normal. \\
\hline
NP & ukuran populasi virtual & Parameter ukuran populasi virtual pada algoritma compact/EDA. \\
\hline
OA & Open Access & Akses terbuka; status lisensi rujukan pustaka. \\
\hline
Permit / deny & izin / tolak & Dua kemungkinan keputusan akses pada log dan kebijakan ABAC. \\
\hline
PMBGA & Probabilistic Model-Building Genetic Algorithm & Nama lain untuk kelas Estimation of Distribution Algorithm. \\
\hline
Policy mining & penambangan kebijakan & Menyimpulkan aturan kebijakan ABAC secara otomatis dari log akses. \\
\hline
Precision & presisi & Proporsi permintaan yang diprediksi izin dan memang benar diizinkan. \\
\hline
Prim (algoritma) & Prim's algorithm & Algoritma pembentukan pohon rentang maksimum yang hemat memori ($O(D)$). \\
\hline
Probability Vector (PV) & vektor probabilitas & Vektor probabilitas marginal per-bit; komponen parameter model CoDA-EDA. \\
\hline
PSO & Particle Swarm Optimization & Optimasi kawanan partikel. \\
\hline
RAG & Retrieval-Augmented Generation & Pembangkitan berbantuan pencarian; disebut pada rujukan pendekatan berbasis LLM. \\
\hline
Rank-biserial & rank-biserial correlation & Ukuran efek (effect size) untuk uji Wilcoxon signed-rank. \\
\hline
RBAC & Role-Based Access Control & Model kontrol akses berbasis peran. \\
\hline
Recall & recall & Proporsi permintaan yang sebenarnya diizinkan dan berhasil tertangkap kebijakan. \\
\hline
Re-mining (cold) & penambangan ulang penuh & Menambang kebijakan dari inisialisasi kosong setelah drift, sebagai pembanding warm-start. \\
\hline
Reservoir sampling & sampling reservoir & Pengambilan subsampel acak seragam hemat memori dengan sekali baca berkas. \\
\hline
RM1--RM4 & Rumusan Masalah 1--4 & Empat rumusan masalah penelitian. \\
\hline
RSS & Resident Set Size & Ukuran memori residen proses; dasar metrik jejak memori. \\
\hline
Separate-and-conquer & sequential covering / pisah-dan-taklukkan & Strategi menambang satu aturan pada satu waktu lalu mengeluarkan permit yang telah tercakup. \\
\hline
SJR & SCImago Journal Rank & Peringkat jurnal (dasar penentuan kuartil Q1--Q4). \\
\hline
Smart Building IoT & --- & Dataset berdomain IoT yang dibangkitkan pihak lain \parencite{diaz-rodriguez_abac_2025}; dipakai untuk validasi eksternal. \\
\hline
SoC & System on Chip & Sistem satu keping (Broadcom BCM2712 pada Raspberry Pi 5). \\
\hline
SP & Special Publication & Publikasi khusus (mis. NIST SP 800-162 tentang ABAC). \\
\hline
Stratified subsample & subsampel terstratifikasi & Subsampel yang mempertahankan rasio permit/deny distribusi aslinya. \\
\hline
T1--T4 & Tujuan 1--4 & Empat tujuan khusus penelitian, selaras satu-satu dengan RM1--RM4. \\
\hline
$\tau$ (tau) & ambang pemicu drift & Ambang divergensi Jensen-Shannon yang memicu pembelajaran ulang struktur dependensi. \\
\hline
TN & True Negative & Permintaan yang benar diklasifikasikan sebagai tolak. \\
\hline
TP & True Positive & Permintaan yang benar diklasifikasikan sebagai izin. \\
\hline
tracemalloc & --- & Pustaka Python pelacak alokasi memori; instrumen metrik memori utama (MemPeak\_alloc). \\
\hline
UBK & Upgraded Black-winged Kite & Algoritma metaheuristik swarm; baseline berpopulasi \parencite{li_attribute_2026}. \\
\hline
UMDA & Univariate Marginal Distribution Algorithm & EDA kanonik berpopulasi-univariat; dipakai sebagai rujukan konseptual EDA (Bab II). \\
\hline
VmHWM & Virtual Memory High Water Mark & Puncak memori tingkat sistem operasi; metrik memori konteks (MemPeak\_sys). \\
\hline
Warm-start & mulai-hangat & Melanjutkan re-optimasi kebijakan dari model sebelumnya alih-alih dari nol setelah drift. \\
\hline
$w_{fp}$, $w_{wsc}$ & bobot fungsi fitness & Bobot penalti false positive ($w_{fp}$) dan kompleksitas struktural ($w_{wsc}$) pada fungsi fitness. \\
\hline
WSC & Weighted Structural Complexity & Kompleksitas struktural kebijakan, yakni jumlah pasangan atribut-nilai pada seluruh aturan. \\
\hline

\end{longtable}
